\scanpage{304}
\mspage{254}
procédant de façon symétrique, ou en factorisant simultanément $p_1$ et $p_2$, et prenant
\diagram{\begin{tikzcd}[ampersand replacement=\&]
\overline X_1\arrow[d]\&\overline X_1\times_S\overline X_2\arrow[l,arrows={-{Bar}}]\arrow[d,dash]\\
S\&\overline X_2\arrow[l]
\end{tikzcd}.}

La loi fonctorielle s'obtient à l'aide du lemme 1 p.~249, qui permet de se borner aux arguments
\diagram{\begin{tikzcd}[ampersand replacement=\&,row sep=small,column sep=small]
X_1\arrow[d,"{p_1}"']\&\\S\&X_2\arrow[l,arrows={-{Bar}},"{p_2}"]
\end{tikzcd}}
tels que $p_1,p_2\in\mathrm{Fib}_W$, de sorte que l'on \unsure{donne|trouve} un foncteur
\diagram{\begin{tikzcd}[ampersand replacement=\&,row sep=small,column sep=small]
X_1\arrow[d]\&\&{\longmapsto}\&X_1\arrow[d]\&X_1\times_S X_2\arrow[l]\arrow[d]\\
S\&X_2\arrow[l]\&\&S\&X_2\arrow[l]
\end{tikzcd}}
compatible aux localiseurs, et passant aux catégories de fractions $\mathrm{Ho}_W(\Psi)$ et $\mathrm{Ho}_W(\Psi^{\sim})$ respectivement.

\uline{NB} En procédant de façon similaire, plus simple (sans besoin d'un recours au lemme 1), on trouve l'existence d'un \uline{adjoint à gauche} $\unsure{i|j}_!$ de $i^*\colon\mathrm{Ho}_W(\Psi^{\sim})\to\mathrm{Ho}_W(\Psi)$, sous réserve de l'existence d'un objet initial $\unsure{\varnothing|\mathscr O|\varphi}$ dans $M$. \unsure{Cet|Ce} foncteur $\unsure{j|i}_!$ est donné par
% archive sheet number: 304
